% -*- mode:LaTex; mode:visual-line; mode:flyspell; fill-column:75-*-

\chapter{Integrated System and Field Validation} \label{chapErwin}

% TODO: Chapter intro — describe the full integrated system (Amiga +
% xArm6, ROS 2, all pipeline components) and the field experiments
% conducted to validate end-to-end performance in a real apple orchard.

\section{Robot Platform} \label{secRobotPlatform}

% TODO: Describe the full physical system.

\subsection{Mobile Base: Amiga} \label{secAmiga}

% TODO: Describe the Farm-ng Amiga robot base.
% - Dimensions, drive configuration (differential/skid steer).
% - Payload capacity, battery life, field navigation capability.
% - Why it was chosen (open platform, orchard-compatible form factor).

\subsection{Manipulator: xArm6} \label{secXArm6}

% TODO: Describe the UFACTORY xArm6 arm.
% - DOF, reach, payload.
% - Mounting configuration on Amiga.
% - Role: position the sensor rig at desired viewpoints.
% - Workspace limitations and collision avoidance.

% \begin{figure}[h]
%   \centering
%   \includegraphics[width=0.75\textwidth]{figs/erwin/robot_platform}
%   \caption{TODO: Photo of the full integrated Amiga + xArm6 robot
%            platform with sensor rig deployed in an orchard row.}
%   \label{figRobotPlatform}
% \end{figure}

\subsection{Sensor Integration and Calibration} \label{secSensorCalib}

% TODO: Describe extrinsic and intrinsic calibration of all cameras
% (FLIR Firefly, Ximea, ZED2i) with respect to the arm end-effector
% and the robot base frame.
% - Calibration target and procedure.
% - Resulting calibration accuracy (reprojection error).
% - TF tree in ROS 2.

\section{System Architecture} \label{secSystemArch}

% TODO: Describe the full software architecture.

\subsection{ROS 2 Integration} \label{secROS2}

% TODO: Describe the ROS 2 node graph.
% - Sensor driver nodes (FLIR, Ximea, ZED2i).
% - Detection inference node (YOLO).
% - Mapping node (semantic OctoMap update).
% - NBV planner node.
% - Arm controller and waypoint execution.
% - Include a node graph / architecture diagram.

% \begin{figure}[h]
%   \centering
%   \includegraphics[width=\textwidth]{figs/erwin/ros2_architecture}
%   \caption{TODO: ROS 2 node graph showing data flow from sensors
%            through detection, mapping, and planning.}
%   \label{figROS2Arch}
% \end{figure}

\subsection{Perception Pipeline} \label{secPerceptionPipeline}

% TODO: Describe end-to-end data flow from image capture to
% semantic map update at runtime.
% - Image capture trigger and synchronization.
% - Inference latency and frame rate.
% - Map update frequency.

\subsection{Planning and Control Loop} \label{secPlanningControl}

% TODO: Describe the closed-loop planning-execution cycle:
% 1. Robot drives to row position.
% 2. NBV planner queries current map; selects next viewpoint.
% 3. Arm moves to viewpoint; images captured.
% 4. Detections update semantic OctoMap.
% 5. Repeat until row complete or budget exhausted.
% Include a flowchart or pseudocode block.

\section{Field Experiments} \label{secFieldExperiments}

% TODO: Describe the real-world experiments.

\subsection{Experimental Site} \label{secFieldSite}

% TODO: Describe the orchard used for field validation.
% - Location (general region, no private address needed).
% - Tree variety, canopy type, row spacing.
% - Disease presence: natural infection or inoculated?
% - Environmental conditions during experiments (dates, weather).

% \begin{figure}[h]
%   \centering
%   \includegraphics[width=\textwidth]{figs/erwin/field_site}
%   \caption{TODO: Photos of the experimental orchard site.}
%   \label{figFieldSite}
% \end{figure}

\subsection{Experimental Protocol} \label{secFieldProtocol}

% TODO: Describe the field experiment procedure.
% - Planner variants tested (Grid Scan, Volumetric NBV, Semantic NBV).
% - Number of trials per planner.
% - Evaluation procedure: how ground truth was established
%   (expert walking survey, Gaussian splat, GPS-tagged disease locations).
% - Metrics collected: detection recall, false positive rate,
%   map coverage, mission time, path length.

\subsection{Field Results} \label{secFieldResults}

% TODO: Report quantitative results from field experiments.

% \begin{table}[h]
%   \centering
%   \caption{TODO: Field experiment results comparing planners.}
%   \label{tabFieldResults}
%   \begin{tabular}{lcccc}
%     \hline
%     Planner & Recall & Precision & Coverage & Time (min) \\
%     \hline
%     Grid Scan & & & & \\
%     Volumetric NBV & & & & \\
%     Semantic NBV (ours) & & & & \\
%     \hline
%   \end{tabular}
% \end{table}

% \begin{figure}[h]
%   \centering
%   \includegraphics[width=\textwidth]{figs/erwin/field_maps}
%   \caption{TODO: Semantic disease maps produced by each planner
%            from the same orchard row.}
%   \label{figFieldMaps}
% \end{figure}

\subsection{Qualitative Observations} \label{secFieldQualitative}

% TODO: Discuss observations from field deployment.
% - Robot behavior in unstructured conditions (terrain, wind, sunlight).
% - Detection performance vs. lab / simulation.
% - Unexpected failure modes encountered and how they were handled.

% \begin{figure}[h]
%   \centering
%   \includegraphics[width=\textwidth]{figs/erwin/field_qualitative}
%   \caption{TODO: Example field detections overlaid on RGB images,
%            and corresponding semantic map updates.}
%   \label{figFieldQualitative}
% \end{figure}

\section{Discussion} \label{secFieldDiscussion}

% TODO: Discuss what the field results tell us.
% - Does semantic NBV outperform baselines in the field as in simulation?
%   If not, explain why (e.g., localization error, detection noise).
% - How does detection performance change from lab to orchard conditions?
% - What are the practical bottlenecks for deployment: compute, sensing,
%   navigation, or planning?
% - What would be needed to scale this to a full commercial orchard?

\section{Summary} \label{secErwinSummary}

% TODO: Summarize:
% - The full integrated Amiga + xArm6 + ROS 2 system.
% - Key field validation results.
% - How simulation results translate (or do not) to real-world performance.
% - Bridge to Chapter~\ref{chapConclusions}: what was demonstrated and
%   what remains to be done.
